\chapter{Terms, Acronyms, and Notation}
\label{chap:notation}

Este capítulo estabelece a terminologia, as siglas e a notação matemática utilizadas ao longo deste trabalho. Para garantir a conformidade com a padronização atual dos protocolos de comunicação, as definições e os símbolos apresentados seguem os documentos oficiais de Request for Comments (RFC) da Internet Engineering Task Force (IETF).

\section{Terms and Definitions}

TBD

\begin{comment}
\todo{FAZER ESSA SEÇÃO}

\begin{description}
    \item[Cryptographic module: ] The set of hardware, software, and/or firmware that implements approved cryptographic functions (including key generation) that are contained within the cryptographic boundary of the module.
    \item[Decapsulation: ] The process of applying the Decaps algorithm of a KEM. This algorithm accepts a KEM ciphertext and the decapsulation key as input and produces a shared secret key as output.
    \item[Encapsulation: ] The process of applying the Encaps algorithm of a KEM. This algorithm accepts the encapsulation key as input, requires private randomness, and produces a shared secret key and an associated ciphertext as output.
    \item[(KEM) Ciphertext: ] A bit string that is produced by encapsulation and used as an input to decapsulation.
    \item[Shared secret key: ] A shared secret that can be used directly as a cryptographic key in symmetric-key cryptography. It does not require additional key derivation.
    \item[Security strength: ] A number associated with the amount of work (i.e., the number of operations) that is required to break a cryptographic algorithm or system.
\end{description}
\end{comment}


% \section{Mathematical Symbols}

\begin{comment}
    
The implementation of ML-KEM relies on polynomial arithmetic over rings. The mathematical notation utilized for the algorithmic description and hardware acceleration is defined below.

\begin{table}[!ht]
\renewcommand{\arraystretch}{1.5}
\centering
\begin{tabular}{@{}lp{13cm}@{}}
\toprule
\textbf{Symbol} & \textbf{Description} \\ \midrule
$n$ & Denotes the integer 256 throughout this document. \\
$q$ & Denotes the prime integer $3329 = 2^8 \cdot 13 + 1$ throughout this document. \\
$\zeta$ & Denotes the integer 17, which is a primitive $n$-th root of unity modulo $q$. \\
$\mathbb{B}$ & The set $\{0, 1, \dots, 255\}$ of unsigned 8-bit integers (bytes). \\
$\mathbb{Z}_q$ & The ring of integers modulo $q$ (i.e., the set $\{0, 1, \dots, q-1\}$ equipped with the operations of addition and multiplication modulo $q$). \\
$R_q$ & The ring $\mathbb{Z}_q[X]/(X^n + 1)$ consisting of polynomials of the form $f = f_0 + f_1X + \dots + f_{255}X^{255}$, where $f_j \in \mathbb{Z}_q$ for all $j$. \\
$T_q$ & The image of $R_q$ under the number-theoretic transform. Its elements are called ``NTT representations'' of polynomials in $R_q$. \\
$\hat{f}$ & The element of $T_q$ that is equal to the NTT representation of a polynomial $f \in R_q$. \\
$f_j$ & The coefficient of $X^j$ of a polynomial $f \in R_q$. \\
$\mathcal{D}_\eta(R_q)$ & A certain distribution of polynomials in $R_q$ with small coefficients, from which noise is sampled. \\
$\times_{T_q}$ & Denotes the operation on coefficient arrays that implements product in the ring $T_q$. \\
$\lceil x \rceil$ & The ceiling of $x$ (i.e., the smallest integer greater than or equal to $x$). \\
$\lfloor x \rfloor$ & The floor of $x$ (i.e., the largest integer less than or equal to $x$). \\
$\lfloor x \rceil$ & The rounding of $x$ to the nearest integer. \\
$A \parallel B$ & The concatenation of two arrays or bit strings $A$ and $B$. \\
$\text{BitRev}_7(r)$ & Bit reversal of a seven-bit integer $r$. \\
\bottomrule
\end{tabular}
\end{table}
\end{comment}